\documentclass[10pt,a4paper]{amsart}
\usepackage[margin=19mm]{geometry}
\usepackage{amsmath,amssymb,mathtools}
\usepackage[scr=rsfs,cal=euler]{mathalfa}
\usepackage{xfrac}
\usepackage[unicode,hidelinks,bookmarks=false]{hyperref}
\newcommand{\ie}{i.e.\ }
\newcommand{\cf}{cf.\ }
\newcommand{\resp}{respectively\ }
\newcommand{\Subsection}{\subsection}
\providecommand{\nobreakdash}{\nobreak\hbox{-}}
\setlength{\emergencystretch}{2em}
\multlinegap0pt
\usepackage{bm}
\usepackage{bbm}
\usepackage{dsfont}
\usepackage{xspace}
\usepackage{cancel}
\usepackage{mathtools}
\usepackage{amsthm}
\makeatletter
\@ifundefined{theorem}{\newtheorem{theorem}{Theorem}}{}
\@ifundefined{lemma}{\newtheorem{lemma}[theorem]{Lemma}}{}
\makeatother
\title[Greedily Constructing Small Quasi-Kernels]{Greedily Constructing Small Quasi-Kernels}
\author{Alexander Clow}
\date{Source: arXiv:2601.11847v1 (CC BY 4.0)}
\begin{document}
\maketitle

\noindent\emph{Source and licence.} Greedily Constructing Small Quasi-Kernels. Version arXiv:2601.11847v1 (CC BY 4.0), distributed under the Creative Commons Attribution 4.0 International licence, as recorded in the arXiv licence metadata for this exact version and shown on the abstract page https://arxiv.org/abs/2601.11847v1. Full rights notice: licenses/arxiv-2601.11847-CC-BY-4.0.txt.\par\smallskip

\noindent\textit{Editorial scope.} Counted unit: the Lemma whose source label is \texttt{Lemma: out-3, min-2} in the article (source0.tex lines 901--908, source label \texttt{Lemma: out-3, min-2}), delivered together with the article's own complete original proof of it (source0.tex lines 910--1038, one \texttt{proof} environment of 129 lines). No mathematics is written, repaired or re-derived: statement and proof are the article's own text line for line. Required context retained uncounted: Lemma: out-3, min-0 (lines 684--691); Lemma: out-3, min-1 (lines 776--783). Every definition, display and label of those lines is retained unchanged. Omitted from this package and not redistributed: 1--683, 692--775, 784--900, 909--909, 1039--1430. Nothing inside a retained excerpt is omitted or altered: every retained body is delivered exactly as the article's lines are written, so no omission receipt of any kind accompanies this package. Citations: the retained excerpts carry no citation command at all, so no bibliography entry of the article is transcribed and no reference is redistributed. Reference adaptation: none was needed and none was made; every reference inside the delivered unit resolves inside it, so no unresolved reference or citation is produced. Printed slips kept literal: no printed word, symbol, hypothesis or number of the counted statement or of its proof was altered. Layout and class: the article's own class is \texttt{article} with packages including authblk, algpseudocode, algorithm, subcaption, amsfonts, graphicx, tabularx, array, color, amsmath, amsthm, amssymb, fullpage, xcolor; the wrapper is the lane's pinned \texttt{amsart} editorial wrapper at 10pt on A4, which restates the theorem-like environments the retained text needs and adds the article's own macro definitions that the retained text needs (none), with extra wrapper packages (bm, bbm, dsfont, xspace, cancel, mathtools). The delivered numbering is the unit's own. Assets: no retained span contains a figure environment, a graphics-inclusion command, a PDF-page inclusion, a table or a TikZ picture; no image file is copied into this package, the package declares no asset dependency and no image file is redistributed with it. Licence: version arXiv:2601.11847v1 of this article is distributed under the Creative Commons Attribution 4.0 International licence, as recorded in the arXiv licence metadata for this exact version and shown on the abstract page https://arxiv.org/abs/2601.11847v1. A full rights notice is delivered at licenses/arxiv-2601.11847-CC-BY-4.0.txt. Characters: every retained excerpt is pure ASCII, so the delivered engine, XeLaTeX with its default Latin Modern text font, reports no missing character.
\begin{lemma}\label{Lemma: out-3, min-0}
   If $D$ is a sourceless oriented graph with maximum out-degree $\Delta^+(D)\leq 3$,
   and no vertex $v$ such that 
    \[
    \frac{4}{3}|N^+(v) \cup N^+(S(v))| \geq |S(v)|+1,
    \]
    then $D$ has minimum out-degree $\delta^+(D)\geq 1$.
\end{lemma}

\begin{lemma}\label{Lemma: out-3, min-1}
   If $D$ is a sourceless oriented graph with maximum out-degree $\Delta^+(D)\leq 3$,
   and no vertex $v$ such that
    \[
    \frac{4}{3}|N^+(v) \cup N^+(S(v))| \geq |S(v)|+1,
    \]
    then $D$ has minimum out-degree $\delta^+(D)\geq 2$.
\end{lemma}

\begin{lemma}\label{Lemma: out-3, min-2}
   If $D$ is a sourceless oriented graph with maximum out-degree $\Delta^+(D)\leq 3$,
   and no vertex $v$ such that
    \[
    \frac{4}{3}|N^+(v) \cup N^+(S(v))| \geq |S(v)|+1,
    \]
    then $D$ has minimum out-degree $\delta^+(D)\geq 3$.
\end{lemma}

\begin{proof}
    Let $D$ be an sourceless oriented graph with maximum out-degree $\Delta^+(D)\leq 3$,
   and no vertex $v$ such that 
    \[
    \frac{4}{3}|N^+(v) \cup N^+(S(v))| \geq |S(v)|+1.
    \]
    By Lemma~\ref{Lemma: out-3, min-1} $D$ has minimum out-degree at least $2$.
    For contradiction suppose that there is a vertex $x$ in $D$ with out-degree $2$.
    Since $D$ is sourceless there exists a vertex $u \in N^-(x)$.
    Let $x,u$ be such vertices, and let $\{y,t\} = N^+(x)$.

    If $\deg^+(u) = 2$, 
    then let $N^+(u) = \{x,w\}$ without loss of generality.
    Observe that $|S(u)|\leq 5$.
    If $S(u) \subseteq N^+(x)$ or $S(u) \subseteq N^+(w)$, then we reach a contradiction by the same argument used in the proof of 
    Lemma~\ref{Lemma: out-3, min-0}.
    Suppose that $S(u) \not\subseteq N^+(x)$ and $S(u) \not\subseteq N^+(w)$.
    Without loss of generality suppose $y \in S(u)$.

    Consider the case where $|N^+(S(u))\setminus N^+(u)]| \geq 3$.
    Then, $|S(u)|\geq 6$, otherwise we can let $v = u$.
    But this is a contradiction, so we suppose $|N^+(S(u))\setminus N^+(u)| \leq 2$.

    Similarly if $|N^+(S(u))\setminus N^+(u)]| = 2$, then $|S(u)| = 5$ so $t \in S(u)$.
    It follows that $S(x)\subseteq N^+(S(u))$.
    As in earlier lemmas $w \notin S(x)$, hence, $|S(x)|\leq 2$.
    Given $\deg^+(x) = 2$ we note that $|S(x)| \geq 2$ otherwise we can take $v =x$, so $|S(x)| = 2$.
    Again since $\deg^+(x) = 2$, and the fact that $D$ is oriented with minimum out-degree $2$, while $S(x)$ is an independent set,
    we observe that $N^+(S(x))\setminus N^+(x) \neq \emptyset$.
    Thus, 
    \[
    |S(x)|+1 = 3 < \frac{4}{3}(3)\leq \frac{4}{3}|N^+(x)\cup N^+(S(x))|
    \]
    letting us take $v = x$, a contradiction.
    We conclude $|N^+(S(u))\setminus N^+(u)]| \leq 1$.
    By the same argument as $x$, $u$ having out-degree $2$ implies $N^+(S(u))\setminus N^+(u) \neq \emptyset$,
    thus we conclude $|N^+(S(u))\setminus N^+(u)]| = 1$.

    Since $|N^+(S(u))\setminus N^+(u)]| = 1$ we note that $|S(u)|\geq 4$, otherwise we can take $v = u$.
    Hence, $N^+(x) \subseteq S(u)$ or $N^+(w) \subseteq S(u)$.
    Let $\Bar{x}$ be this vertex. Then as should be familiar $S(\Bar{x}) \cap N^+(u) = \emptyset$.
    Since $N^+(\Bar{x}) \subseteq S(u)$, we note $S(\Bar{x}) \subseteq N^+(S(u))$,
    so $|S(\Bar{x})|\leq 1$.
    But $\deg^+(\Bar{x})\geq 2$ since $D$ has minimum out-degree $2$.
    Thus, we can take $v = \Bar{x}$ a contradiction.
    We conclude that $\deg^+(u) = 3$.


    Without loss of generality let $N^+(u) = \{x,w,z\}$,
    and without loss of generality suppose that $|N^+(w)\cap S(u)|\geq |N^+(z)\cap S(u)|$.
    Then $|S(u)|\leq 8$.
    If $|N^+(S(u))\setminus N^+(u)|\geq 4$, then
    \[
    |S(u)|+1 \leq 9 < \frac{4}{3}(7) \leq \frac{4}{3}|N^+(u) \cup N^+(S(u))|
    \]
    letting us take $v = u$, so $|N^+(S(u))\setminus N^+(u)|\leq 3$.
    Consider the case where $N^+(x)\subseteq S(u)$.
    Then,  $S(x) \subseteq N^+(S(u))\setminus N^+(u)$ implying $|S(x)|\leq 3$.
    Since $x$ has degree $2$, and $D$ has minimum degree $2$, $N^+(S(x))\setminus N^+(x) \neq \emptyset$.
    Hence, 
    $|S(x)|+1 \leq 4 \leq \frac{4}{3}(3) = \frac{4}{3}|N^+(x) \cup N^+(S(x))|$
    allowing us to take $v = x$ a contradiction.
    We conclude $N^+(x)\subseteq S(u)$, which implies $|S(u)|\leq 7$.
    Since $|S(u)|\leq 7$ we note that $|N^+(S(u))\setminus N^+(u)|= 3$ implies 
    \[
    |S(u)|+1 \leq 8 = \frac{4}{3}(6) \leq \frac{4}{3}|N^+(u) \cup N^+(S(u))|
    \]
    letting us take $v = u$, so $|N^+(S(u))\setminus N^+(u)|\leq 2$.
    Finally, observe that since $N^+(x)\not\subset S(u)$, 
    if $|N^+(S(u))\setminus N^+(u)|= 2$, implies  $|N^+(w)\cap S(u)| = 3$.

    
    Consider the case where $N^+(w)\subseteq S(u)$.
    As earlier, $x,z \notin S(w)$, so $S(w) \cap N^+(u) = \emptyset$.
    Since, $N^+(w)\subseteq S(u)$, we note that $S(w)\subseteq N^+(S(u))$.
    If $\deg^+(w) = 2$, then 
    \[
    2 \leq |S(w)| \leq |N^+(S(u))\setminus N^+(u)|
    \]
    otherwise we can choose $v = w$.
    But this is a contradiction, because we have shown $|N^+(S(u))\setminus N^+(u)|\geq 2$ forces 
    $\deg^+(w) = |N^+(w)\cap S(u)| = 3> 2 = \deg^+(w)$.
    So we suppose that $\deg^+(w) = 3$.
    As before, $S(w)\cap N^+(u) = \emptyset$ and $S(w) \subseteq N^+(S(u))$.
    Since $\deg^+(w) = 3$ we note that $|S(w)|\geq 4$,
    otherwise we can take $v = w$,
    implying that
    $4 \leq |S(w)| \leq |N^+(S(u))\setminus N^+(u)|$ given $S(w) \cap N^+(u) = \emptyset$.
    But this implies 
    \[
    |S(u)| + 1 \leq 8 < \frac{4}{3}(7) \leq \frac{4}{3}|N^+(u) \cup N^+(S(u))|.
    \]
    Since this would let us take $v = u$, we conclude that $N^+(w)\not\subseteq S(u)$.

    
    Since, $N^+(x)\not\subseteq S(u)$, and 
    $|N^+(w)\cap S(u)|\geq |N^+(z)\cap S(u)|$ we note that if $|N^+(w)\cap S(u)|\leq 1$, then $|S(u)|\leq 3$.
    This would allow us to take $v = u$ since $\deg^+(u) = 3$,
    so we suppose $|N^+(w)\cap S(u)|\geq 2$.
    From here $N^+(w)\not\subseteq S(u)$ implies that $\deg^+(w) = 3$ and $|N^+(w)\cap S(u)|= 2$.
    Now, $|N^+(w)\cap S(u)|= 2$, $N^+(x)\not\subseteq S(u)$, and 
    $|N^+(w)\cap S(u)|\geq |N^+(z)\cap S(u)|$ implies $|S(u)|\leq 5$.

    As before, $S(w) \cap N^+(u) = \emptyset$.
    This implies $|S(w)| - 3 \leq |N^+(S(u))|$, given the unique vertex in $N^+(w)\setminus S(u)$
    contributes at most $3$ vertices to $S(w)$.
    Since $\deg^+(w) = 3$ we note that $|S(w)|\geq 4$,
    otherwise we can take $v = w$,
    implying that $1 \leq |S(w)| - 3 \leq |N^+(S(u))\setminus N^+(u)|$.
    If $|S(u)|\leq 4$, then 
    \[
    5 = |S(u)| + 1 < \frac{4}{3}(4) \leq \frac{4}{3}|N^+(u) \cup N^+(S(u))|.
    \]
    allowing us to take $v = u$.
    Hence, we conclude that $|S(u)| = 5$ which forces $N^+(x)\cap S(u) \neq \emptyset$.
    Since $|S(u)|\leq 5$ we note that $|N^+(S(u))\setminus N^+(u)| = 1$, otherwise we can take $v = u$.
    Without loss of generality suppose that $y \in S(u)$ and $t\notin S(u)$.

    As before, $S(x) \cap N^+(u) = \emptyset$.
    This implies $|S(x)| - 3 \leq |N^+(S(u))\setminus N^+(u)| = 1$, given the unique vertex in $N^+(x)\setminus S(u)$, so $|S(x)|\leq 4$.
    Since $\deg^+(x) = 2$ and $D$ has minimum out-degree $2$, we note that $|N^+(S(x))\setminus N^+(x)|\geq 1$.
    It follows that $|S(x)| \geq 4$, implying $|S(x)|= 4$, which forces
    $|N^+(S(x))\setminus N^+(x)|= 1$ and $|N^+(t) \cap S(x)| = 3$.
    Since $|N^+(t) \cap S(x)| = 3$, $N^+(t) \subseteq S(x)$ and $\deg^+(t) = 3$.
    But $\deg^+(t) = 3$ implies $|S(t)|\geq 4$ while $N^+(t) \subseteq S(x)$ implies
    $S(t) \subseteq N^+(S(x))\setminus N^+(x)$.
    This implies $4 \leq |S(t)| \leq |N^+(S(x))\setminus N^+(x)| =1$ a contradiction.
    This completes the proof.
\end{proof}

\end{document}
